\documentclass[conference]{IEEEtran}
\IEEEoverridecommandlockouts
% The preceding line is only needed to identify funding in the first footnote. If that is unneeded, please comment it out.
\usepackage{cite}
\usepackage{url}
\usepackage{etoolbox}
\usepackage{amsmath,amssymb,amsfonts}
\usepackage{algorithmic}
\usepackage{graphicx}
\usepackage{textcomp}
\usepackage{xcolor}
\usepackage{multirow}
\usepackage{booktabs}
\usepackage[utf8]{inputenc}
\usepackage{newunicodechar}
\newunicodechar{−}{-}
\def\BibTeX{{\rm B\kern-.05em{\sc i\kern-.025em b}\kern-.08em
    T\kern-.1667em\lower.7ex\hbox{E}\kern-.125emX}}

% Float and spacing control to guarantee strictly 4 pages body + 1 page references
\setcounter{dbltopnumber}{2}
\renewcommand{\dbltopfraction}{0.9}
\renewcommand{\topfraction}{0.9}
\renewcommand{\textfraction}{0.1}
\renewcommand{\floatpagefraction}{0.8}
\setlength{\textfloatsep}{5pt plus 1pt minus 1pt}
\setlength{\floatsep}{5pt plus 1pt minus 1pt}
\setlength{\intextsep}{5pt plus 1pt minus 1pt}
\setlength{\dbltextfloatsep}{5pt plus 1pt minus 1pt}
\setlength{\dblfloatsep}{5pt plus 1pt minus 1pt}

\begin{document}

\title{Decoding Mind Wandering Through Neural Dynamics and Complementary Features\vspace{-4mm}
}

\author{
\IEEEauthorblockN{
Shubham\textsuperscript{1},
Rishwanth JVK\textsuperscript{2},
Amita Giri\textsuperscript{1}
}

\IEEEauthorblockA{
\textsuperscript{1}\textit{Department of Electronics and Communication Engineering, Indian Institute of Technology, Roorkee, 247667, India}\\
\textsuperscript{2}\textit{Department of Computer Science and Engineering, Manipal Institute of Technology, Karnataka,576104,India}\\
Emails: shubham1,amita.giri@ece.iitr.ac.in, rishwanthjvk06@gmail.com
}
}

\maketitle

\begin{abstract}
\textcolor{blue}{Mind Wandering} is a common cognitive state characterized by a shift of attention from ongoing tasks toward internally generated thoughts. Objective assessment using electroencephalography (EEG) remains challenging due to the complex, multiscale nature of brain activity. We propose an integrated multi-domain EEG framework for characterizing and classifying \textcolor{blue}{Mind Wandering} across two complementary datasets with distinct channel densities and experimental paradigms. The framework integrates spectral, temporal, complexity, entropy, amplitude-envelope, and functional connectivity features. On pooled epoch data, \textcolor{blue}{simple linear models proved good enough to effectively classify Focus and Mind Wandering states, achieving up to $88.06\%$ accuracy on high-density EEG and performing on par with a deep multilayer perceptron ($86.06\%$), demonstrating that classifier complexity is secondary to rich multi-domain feature representation}. \textcolor{blue}{Beyond conventional $\alpha$ and $\theta$ rhythms, the analysis revealed robust multi-band biomarkers, including near-identical Delta baseline convergence ($\sim+0.082$ vs $+0.083\text{ log}_{10}(\mu\text{V}^2/\text{Hz})$), progressive Delta drift during Mind Wandering, steep within-state Beta trajectory divergence ($p < 10^{-80}$), and Gamma sensory decoupling suppression. Furthermore, Hjorth parameters demonstrated the highest per-feature information density across both datasets, while posterior scalp regions dominated standalone spatial decoding ($73.83\%$ accuracy)}. These unified findings establish multi-domain EEG dynamics as a powerful tool for objective \textcolor{blue}{Mind Wandering} characterization.
\end{abstract}

\textbf{\textit{Keywords---}\textit{\textcolor{blue}{Mind Wandering}}, \textit{EEG}, \textit{Spectral Features}, \textit{Nonlinear Features}, \textit{Cognitive State Classification}.}

\section{Introduction}

\textcolor{blue}{Mind Wandering (MW)} is an involuntary attentional shift from external tasks to internal mental processing \cite{smallwood2013restless}, occurring across daily activities and leading to task lapses \cite{killingsworth2010wandering}. Neuroimaging links internally directed cognition to default mode network (DMN) recruitment \cite{raichle2001default} and dynamic executive control interactions \cite{christoff2016mind}. Meditation paradigms provide an ideal experimental framework for studying attentional shifts, offering structured transitions between sustained task engagement and spontaneous cognition \cite{lutz2008attention,tang2015neuroscience,cahn2013meditation}. Electroencephalography (EEG) offers millisecond temporal resolution to capture these transitional neural dynamics \cite{KAM2022119372}. Experience-sampling studies have identified spectral shifts during transitions between \textcolor{blue}{Focus and Mind Wandering (MW)}.

While literature predominantly emphasizes $\alpha$ and $\theta$ oscillations \cite{KAM2022119372}, \textcolor{blue}{relying solely on these rhythms overlooks critical neural dynamics: slow-wave $\delta$ power indexes arousal and vigilance degradation, $\gamma$ suppression marks sensory decoupling from external stimuli, and $\beta$ oscillations capture sensorimotor maintenance and motor cessation.} Beyond spectral measures, nonlinear complexity and entropy quantify broader shifts in cortical organization \cite{lu2022nonlinear,berkovich2012mindfulness}. Machine-learning methods have automated \textcolor{blue}{Mind Wandering} detection via spectral features and classifiers \cite{jin2020distinguishing,compton2019wandering,hosseini2019deep}. However, existing approaches often rely on predefined bands or isolated feature types without systematically examining multi-band temporal trajectories or feature information density across recording montages.

In this study, we investigate EEG biomarkers and classification across two publicly available datasets. The main contributions are: (1) an integrated multi-domain EEG representation combining spectral, temporal, complexity, envelope, and connectivity metrics; (2) \textcolor{blue}{identification of consistent physiological patterns across frequency bands: convergent Delta baseline elevation ($\sim+0.082$ to $+0.083\text{ log}_{10}(\mu\text{V}^2/\text{Hz})$), progressive Delta drift during Mind Wandering, robust Theta elevation, Gamma suppression indexing sensory decoupling, and steep within-state Beta divergence ($p < 10^{-80}$); (3) feature-importance analysis revealing the foundation of band power ($49.9\%$ in Dataset~1, $33.7\%$ in Dataset~2) and envelope dynamics ($29.5\%$), paired with the highest information density in Hjorth parameters across both datasets ($20.8\%$ and $20.1\%$ total), while phase-based functional connectivity is redundant ($<6.8\%$ in Dataset~1, $0.00\%$ in Dataset~2); and (4) classification demonstrating that simple, lightweight linear models are good enough to reliably classify Focus and Mind Wandering states without requiring complex deep networks, reaching $88.06\%$ accuracy with Logistic Regression on high-density EEG, while posterior scalp electrodes dominate standalone spatial decoding ($73.83\%$)}.



\section{Proposed Methodology}

\subsection{Datasets and Standardized Preprocessing}

We utilized two publicly available EEG datasets \cite{grandchamp2014oculometric,rodlarios}. Dataset~1 \cite{rodlarios} comprises 19 EEG channels (Nexus-32 system, 10--20 montage) sampled at 512~Hz with VEOG/HEOG at 2048~Hz\textcolor{blue}{, from 25 retained participants (40 paired epochs each, yielding 1,000 balanced Focus and Mind Wandering epochs)}. Dataset~2 \cite{grandchamp2014oculometric} comprises 64-channel BIOSEMI recordings (1024~Hz) from healthy individuals performing a reverse breath-counting task with probe-prompted button presses indicating \textcolor{blue}{Mind Wandering (MW)} and refocusing events, yielding \textcolor{blue}{1,047 epochs (472 Mind Wandering and 575 Focus)}.

To ensure rigorous cross-dataset comparability, \textcolor{blue}{preprocessing was standardized: both datasets were resampled to 256~Hz and band-pass filtered between 0.5--45~Hz}. For Dataset~1, channels were converted to volts, spherical-spline interpolated for bad channels, and average-referenced. \textcolor{blue}{Focus and Mind Wandering} epochs were extracted from $-5$ to $-0.1$~s and edge-padded to 5~s. Ocular artifacts were eliminated using EOG-assisted ICA ($|r|\geq 0.30$) \cite{delorme2004eeglab}. For Dataset~2, recordings were evaluated under \textcolor{blue}{Raw and Artifact removed} conditions (ICA components with $|r|\geq 0.30$ relative to auxiliary EXG channels rejected). \textcolor{blue}{Mind Wandering} epochs ($-5$ to $-0.1$~s) and \textcolor{blue}{Focus} epochs ($+1$ to $+6$~s) were edge-padded to uniform 5-s windows. Matching pipelines were applied to both branches to assess \textcolor{blue}{artifact removal}.

\subsection{Multi-Domain Feature Extraction Framework}

Each 5-s epoch was partitioned into 1.5-s windows with a 0.5-s step. Features were computed across five bands: $\delta$ (1--4~Hz), $\theta$ (4--8~Hz), $\alpha$ (8--12~Hz), \textcolor{blue}{$\beta$ (12--30~Hz)}, and $\gamma$ (30--45~Hz).
The framework integrates multiple domains: statistical \& time-domain (mean, variance, standard deviation, RMS, kurtosis, skewness, peak-to-peak amplitude, energy); Hjorth dynamics (activity, mobility, complexity); entropy \& complexity (Shannon, differential, spectral entropy) \cite{zhang2018integration}; spectral features (Welch PSD-derived absolute and relative band power, total power, dominant frequency) \cite{safi2021early}; and envelope \& synchrony (Hilbert amplitude envelopes for $\alpha$ and $\theta$ bands, extracting mean, variance, CV, burst count with threshold $\mu+2\sigma$, pairwise PLV, and $\alpha$-envelope inter-channel synchronization) \cite{o2017detecting}, preserving channel identity and spatial topography.
The window feature representation is $\mathbf{f}^{(w)} = [f_1^{(w)}, f_2^{(w)}, \ldots, f_D^{(w)}]^{T}$, where $D$ is feature dimensionality \textcolor{blue}{($D = 665$ for Dataset~1; $D = 4,096+$ for Dataset~2 with pairwise PLV)}. Inter-channel $\alpha$ synchronization for channel $i$ is:
\begin{equation}
S_i^{\alpha} = \frac{1}{C-1} \sum_{j=1, j\neq i}^{C} \mathrm{corr}\left(A_{i,\alpha}(t), A_{j,\alpha}(t)\right),
\end{equation}
where $A_{i,\alpha}(t)$ denotes the $\alpha$ amplitude envelope across $C$ channels.

\subsection{Feature Importance, Ablation, and Temporal Dynamics}

Extra Trees classifiers \cite{geurts2006extremely} were trained on \textcolor{blue}{Raw and Artifact removed} EEG under group-wise cross-validation to compute band importance $I_b^{\mathrm{total}} = \sum_{j\in\mathcal{J}_b} I_j$\textcolor{blue}{, where $I_j$ is cross-validated Gini Mean Decrease in Impurity (MDI)}. Systematic ablation evaluated feature families, scalp regions (anterior, central--temporal, posterior), and frequency bands via omission and isolation under grouped validation with SVM and Logistic Regression (LR). \textcolor{blue}{For within-epoch dynamics, per-epoch trajectory slopes were fitted across 1-s sliding windows (0.5-s step); population trends were tested via OLS with Huber-White cluster-robust standard errors grouped by subject, assessing state slopes via Wald $z$-tests ($H_0: \text{slope}=0$, $\alpha=0.05$)}.

%%%%%%%%%%%%%%%%%%%%%%%%
\begin{figure*}[t]
    \centering
    \includegraphics[width=\linewidth]{1.png}
    \caption{Feature importance per frequency band for (a) Dataset~1 and
    (b) Dataset~2.}
    \label{fig:band_importance}
\end{figure*}
\begin{figure*}[h]
    \centering
    \includegraphics[width=\textwidth]{2.png}
    \caption{Power trajectories of EEG bands for (a) Dataset~1 and
    (b) Dataset~2, and Delta-band power over time for (c) Dataset~1 and
    (d) Dataset~2\textcolor{blue}{; shaded error bands denote $\pm 1$ SEM}.}
    \label{fig:power_trajectories}
\end{figure*}
\begin{figure*}[h]
    \centering
    \includegraphics[width=\linewidth]{3.png}
    \caption{Feature family comparison for (a) Dataset~1 and (b) Dataset~2.}
    \label{fig:feature_family}
\end{figure*}

\subsection{Classification and Evaluation Pipeline}

\textcolor{blue}{To evaluate whether simple, standard classifiers are good enough to distinguish cognitive states rather than conducting an exhaustive algorithmic benchmark, lightweight linear and kernel models (Logistic Regression and SVM with an RBF kernel) served as the primary classification pipeline. A deep multilayer perceptron (DeepMLP: $D \rightarrow 1024 \rightarrow 512 \rightarrow 256 \rightarrow 128 \rightarrow 1$ with batch normalization, GELU, dropout $0.50$--$0.20$, and AdamW) was included as a high-capacity reference model under Raw and Artifact removed conditions. For SVM and LR, 5-fold stratified cross-validation on integrated 5.0-s epoch features assessed global decision boundaries, while DeepMLP was trained on sliding 1.5-s window vectors with foldwise scaling, early stopping, and out-of-fold window-to-epoch probability averaging ($P = \frac{1}{W}\sum_{w=1}^W P_w \ge 0.5$). Evaluation metrics included accuracy, precision, recall, F1-score, and ROC-AUC, computed on out-of-fold predictions.}



\section{Results and Discussion}

\subsection{Multi-Band Spectral Correlates and Temporal Dynamics}

Band-wise feature importance (Fig.~\ref{fig:band_importance}) demonstrated distinct spectral signatures. In Dataset~1, $\alpha$ ($33.50\%$) and $\theta$ ($31.85\%$) dominated, followed by $\gamma$ ($11.87\%$), $\beta$ ($11.47\%$), and $\delta$ ($11.32\%$). In Dataset~2, $\alpha$ carried the highest importance ($40.94\%$ Raw, $39.26\%$ \textcolor{blue}{Artifact removed}), followed by $\theta$ ($25.62\%$ Raw, $26.07\%$ \textcolor{blue}{Artifact removed}), $\beta$ ($14.91\%$--$15.34\%$), $\gamma$ ($10.69\%$--$11.10\%$), and $\delta$ ($7.83\%$--$8.22\%$).

Consistent with these distributions, $\theta$ power increased during \textcolor{blue}{Mind Wandering} relative to \textcolor{blue}{Focus} in both datasets, reaching high significance in Dataset~2 \textcolor{blue}{($\hat{\beta} = +0.0346$, $z = 6.4941$, $p = 8.36\times 10^{-11}$)} and a positive trend in Dataset~1 \textcolor{blue}{($\hat{\beta} = +0.0774$, $p = 0.1455$)}. Furthermore, Dataset~1 within-epoch analysis revealed that \textcolor{blue}{$\theta$ power significantly increased over time during Mind Wandering ($\text{slope} = +0.0080$, $p = 0.0195$)}, indicating progressive immersion in internal cognition.

\textcolor{blue}{Crucially, mean $\delta$ power exhibited near-identical positive elevation during Mind Wandering across datasets ($\hat{\beta} = +0.0823\text{ log}_{10}(\mu\text{V}^2/\text{Hz})$ in Dataset~1; $\hat{\beta} = +0.0833\text{ log}_{10}(\mu\text{V}^2/\text{Hz})$ in Dataset~2), demonstrating a remarkably stable baseline shift toward lower physiological vigilance.} Temporal modeling (Fig.~\ref{fig:power_trajectories}) revealed that \textcolor{blue}{$\delta$ power rose over time during Mind Wandering in both datasets ($\text{slope} = +0.0041\text{ s}^{-1}$ in Dataset~1; $\text{slope} = +0.0182\text{ s}^{-1}$, $p < 0.001$ in Dataset~2), capturing cumulative vigilance decline, whereas during Focus, $\delta$ power decreased in Dataset~2 ($\text{slope} = -0.0069\text{ s}^{-1}$, $p = 0.0589$), indicating maintained alertness, while remaining stable in Dataset~1 ($\text{slope} = +0.0037\text{ s}^{-1}$)}.

\textcolor{blue}{Beyond low frequencies, $\gamma$ power (30--45~Hz) exhibited consistent suppression during Mind Wandering across datasets ($\hat{\beta} = -0.0183$ in Dataset~1; $\hat{\beta} = -0.0329$, $p = 0.0972$ in Dataset~2), accompanied by downward Focus trajectories ($\text{slope} = -0.0004\text{ s}^{-1}$ in Dataset~1; $\text{slope} = -0.0089\text{ s}^{-1}$ in Dataset~2), corroborating sensory decoupling from external stimuli. Finally, while mean $\beta$ differences (12--30~Hz) were non-significant ($\hat{\beta} = -0.0661$ in Dataset~1; $\hat{\beta} = -0.0051$ in Dataset~2), temporal trajectory analysis in Dataset~2 showed dramatic within-state divergence: $\beta$ dropped sharply during Focus ($\text{slope} = -0.0227\text{ s}^{-1}$, $z = -19.92$, $p = 2.87\times 10^{-88}$) but increased during Mind Wandering ($\text{slope} = +0.0088\text{ s}^{-1}$, $p = 0.0491$), establishing temporal $\beta$ modulation as an active marker of task engagement.}

\subsection{Classification Performance and Feature Information Density}

Classification results are summarized in Table~\ref{tab:classification_comparison}. \textcolor{blue}{Crucially, the classification results demonstrate that simple, standard models are entirely good enough to classify between Focus and Mind Wandering states, obviating the need for complex deep learning architectures. Rather than presenting a competitive benchmark between modeling algorithms, our analysis reveals that basic classifiers reliably capture the cognitive boundary once supplied with multi-domain EEG features. On pooled data in Dataset~2, a standard regularized Logistic Regression model achieved the highest overall accuracy of $88.06\%$ (F1: $86.60\%$, ROC-AUC: $94.28\%$) on Raw data and $86.06\%$ on Artifact removed data, matching or slightly exceeding both the DeepMLP ($85.96\%$ Raw / $86.06\%$ Artifact removed) and SVM ($84.91\%$). In Dataset~1, standard SVM ($68.30\%$ Raw / $68.00\%$ Artifact removed) likewise performed on par with the DeepMLP ($67.40\%$ Raw / $68.00\%$ Artifact removed). These findings confirm that classification effectiveness is driven by comprehensive feature engineering—such as spectral power, envelope dynamics, and Hjorth parameters—rather than classifier complexity. Furthermore, artifact removal did not degrade classification fidelity across models (e.g., SVM yielded an identical $84.91\%$ across Raw and Artifact removed conditions in Dataset~2), confirming robust representation stability.}

\begin{table}[h]
\centering
\caption{Classification performance on Dataset~1 and Dataset~2.}
\label{tab:classification_comparison}
\resizebox{\linewidth}{!}{%
\begin{tabular}{lllccccc}
\toprule
Dataset & Condition & Model & Accuracy & Precision & Recall & F1 & ROC-AUC \\
\midrule
\multirow{6}{*}{Dataset 1}
& \multirow{3}{*}{Raw}
& SVM & 68.30 & 65.21 & 58.16 & 61.48 & 72.19 \\
&& LR & 63.20 & 57.24 & 60.92 & 59.02 & 65.61 \\
&& Deep MLP & 67.40 & 63.39 & 59.31 & 61.28 & 71.41 \\
\cline{2-8}
& \multirow{3}{*}{\textcolor{blue}{Artifact removed}}
& SVM & 68.00 & 65.84 & 54.94 & 59.90 & 72.10 \\
&& LR & 61.10 & 55.02 & 57.93 & 56.44 & 65.65 \\
&& Deep MLP & 68.00 & 65.67 & 55.40 & 60.10 & 71.69 \\
\midrule
\multirow{6}{*}{Dataset 2}
& \multirow{3}{*}{Raw}
& SVM & 84.91 & 87.03 & 78.18 & 82.37 & 93.30 \\
&& LR & 88.06 & 87.64 & 85.59 & 86.60 & 94.28 \\
&& Deep MLP & 85.96 & 85.10 & 83.47 & 84.28 & 93.69 \\
\cline{2-8}
& \multirow{3}{*}{\textcolor{blue}{Artifact removed}}
& SVM & 84.91 & 86.85 & 78.39 & 82.41 & 92.80 \\
&& LR & 86.06 & 85.75 & 82.84 & 84.27 & 92.70 \\
&& Deep MLP & 86.06 & 84.39 & 84.75 & 84.57 & 93.10 \\
\bottomrule
\end{tabular}%
}
\end{table}

Feature-family analysis (Fig.~\ref{fig:feature_family}) demonstrated that \textcolor{blue}{band power provided the largest aggregate foundation ($49.93\%$ in Dataset~1, $33.66\%$ in Dataset~2, supplemented by $29.47\%$ envelope dynamics). Crucially, Hjorth parameters (Activity, Mobility, Complexity) carried the highest information density across both datasets, accounting for $20.85\%$ total importance across 38 features in Dataset~1 ($0.005486$ per feature) and surging to $20.05\%$ across 128 features in Dataset~2 ($0.001567$ per feature). Conversely, functional connectivity (PLV) contributed minimally in Dataset~1 ($6.76\%$) and zero in Dataset~2 ($0.00\%$), despite constituting $76\%$ of extracted features (4,096 features), proving that Mind Wandering is indexed by localized spectral and waveform dynamics rather than remote phase synchrony.}



\subsection{Feature Domain and Spatial Ablation Analysis}

Ablation results in Table~\ref{tab:ablation_comparison} highlight domain and spatial dependencies. In Dataset~1, $\theta$ was the strongest isolated band \textcolor{blue}{(Acc.~$=59.67\%$, F1~$=51.27\%$)}, and central--temporal sensors provided the best isolated spatial representation \textcolor{blue}{(Acc.~$=62.70\%$, F1~$=56.68\%$, AUC~$=66.34\%$, exceeding the full montage baseline of $61.50\%$)}. In Dataset~2, $\alpha$ was the strongest isolated band \textcolor{blue}{(Acc.~$=62.45\%$, F1~$=56.97\%$, AUC~$=65.30\%$)}, and omitting $\alpha$ induced \textcolor{blue}{the steepest frequency degradation ($65.36\%$ to $60.20\%$)}.

\textcolor{blue}{Crucially, in the spatial domain for Dataset~2, posterior electrodes produced the highest isolated accuracy ($73.83\%$, F1~$=70.73\%$, AUC~$=78.38\%$), markedly outperforming frontal ($62.85\%$) and central--temporal ($65.71\%$) subsets, and even surpassing the full 64-channel baseline ($68.39\%$). This establishes posterior sensory gating as a paramount standalone hub of Mind Wandering decoding in dense arrays.}

\begin{table}[h]
\centering
\caption{Ablation analysis using SVM on raw EEG \textcolor{blue}{(corresponding LR trends summarized in text)}.}
\label{tab:ablation_comparison}
\resizebox{\linewidth}{!}{%
\begin{tabular}{llrrr|rrr}
\toprule
& & \multicolumn{3}{c|}{Dataset 1} & \multicolumn{3}{c}{Dataset 2} \\
Category & Variant & Acc. & F1 & AUC & Acc. & F1 & AUC \\
\midrule
Baseline
& Full feature set
& 61.50 & 55.70 & 65.37
& 68.39 & 59.39 & 77.53 \\
\midrule
\multirow{2}{*}{Feature family}
& Best omission (Omit Entropy / Omit Hjorth)
& 61.80 & 55.89 & 65.27
& 68.77 & 61.12 & 76.72 \\
& Most damaging omission (Omit Spectral)
& 59.70 & 54.87 & 63.30
& 60.94 & 53.26 & 70.31 \\
\midrule
\multirow{2}{*}{Feature family}
& Best isolated (Only Hjorth / Only Spectral)
& 59.10 & 55.30 & 62.38
& 65.52 & 57.78 & 75.77 \\
& Most damaging isolated (Only Entropy / Only Hjorth)
& 56.60 & 54.51 & 60.61
& 52.63 & 37.22 & 56.27 \\
\midrule
\multirow{2}{*}{Scalp region}
& Best omission (Omit Frontal)
& 63.10 & 56.64 & 65.55
& 67.72 & 55.53 & 78.55 \\
& Most damaging omission (Omit Central/Temporal / Omit Posterior)
& 60.10 & 54.71 & 62.52
& 66.19 & 55.75 & 72.93 \\
\midrule
\multirow{2}{*}{Scalp region}
& Best isolated (Only Central/Temporal / Only Posterior)
& 62.70 & 56.68 & 66.34
& 73.83 & 70.73 & 78.38 \\
& Most damaging isolated (Only Frontal)
& 58.60 & 54.51 & 62.27
& 62.85 & 55.13 & 70.86 \\
\midrule
\multirow{2}{*}{Frequency band}
& Best omission (Omit Alpha / Omit Gamma)
& 58.72 & 52.79 & 61.23
& 65.89 & 60.58 & 69.37 \\
& Most damaging omission (Omit Gamma / Omit Alpha)
& 57.88 & 51.97 & 60.37
& 60.20 & 54.17 & 63.54 \\
\midrule
\multirow{2}{*}{Frequency band}
& Best isolated (Only Theta / Only Alpha)
& 59.67 & 51.27 & 59.50
& 62.45 & 56.97 & 65.30 \\
& Most damaging isolated (Only Delta / Only Gamma)
& 53.82 & 50.30 & 56.03
& 53.25 & 36.78 & 50.11 \\
\bottomrule
\end{tabular}%
}
\end{table}

\section{Conclusion}

This study presented an integrated multi-domain EEG framework for characterizing \textcolor{blue}{Mind Wandering} across two diverse recording setups. The analysis uncovered \textcolor{blue}{robust multi-band biomarkers spanning slow, intermediate, and fast oscillations: convergent Delta baseline elevation ($\sim+0.082$ to $+0.083\text{ log}_{10}(\mu\text{V}^2/\text{Hz})$), progressive Delta drift during Mind Wandering, robust Theta elevation, Gamma sensory decoupling suppression, and steep within-state Beta divergence ($p < 10^{-80}$)}. Feature analysis established that \textcolor{blue}{band power and envelope features provide the foundation, while Hjorth parameters achieve the highest information density across both datasets, contrasted with functional connectivity redundancy. Classification analysis demonstrated that simple models are entirely good enough to classify between Focus and Mind Wandering states (reaching up to $88.06\%$ accuracy with Logistic Regression), performing on par with deep multilayer architectures and confirming that discriminative power stems primarily from comprehensive multi-domain feature design rather than model complexity. Furthermore, artifact removal preserved discriminative fidelity across conditions, and posterior electrodes dominated standalone spatial decoding ($73.83\%$)}. These findings demonstrate that integrated multi-domain EEG representations effectively index transitional dynamics between \textcolor{blue}{Focus and internally directed Mind Wandering states}.

\section*{Acknowledgment}
This research was supported by the Anusandhan National Research Foundation (ANRF), Department of Science and Technology, Government of India, under grant number ANRF/ECRG/2024/004043/ENS.

\clearpage
\bibliographystyle{IEEEtran}
\bibliography{references}

\end{document}
